% =========================================================
\chapter{Les céramiques}
% =========================================================

\begin{rembox}[title=En bref]
\textbf{De quoi parle ce chapitre ?} Première grande classe de matériaux : les céramiques
(pierres, briques et porcelaines, verres, ciments et bétons, céramiques techniques). On voit
leurs propriétés (rigides, dures, réfractaires\ldots{} mais \textbf{fragiles}), leur structure
atomique (cristaux ioniques ou covalents, verres amorphes), leur élaboration (argiles,
frittage, mise en forme du verre, contraintes thermiques) et surtout leur
\textbf{comportement mécanique} : rupture fragile, \textbf{statistique de Weibull}, rupture
en compression, fissuration lente. \\
\textbf{Le fil rouge} : tout découle de $K = Y\sigma\sqrt{\pi a}$ (chapitre~1) appliqué à des
matériaux sans plasticité et pleins de petits défauts.
\end{rembox}

\section{Panorama : une famille rigide, dure\ldots{} et fragile}

\begin{defnbox}[title=Les grandes familles de céramiques]
\begin{enumerate}
  \item \textbf{Naturelles} : calcaire/marbre, grès, granite, glace.
  \item \textbf{Poreuses et vitrifiées} : porcelaine, faïence, briques (à base d'argiles).
  \item \textbf{Verres} : verre sodocalcique ($\approx$ 70\,\% SiO$_2$, 10\,\% CaO, 15\,\%
        Na$_2$O), borosilicaté (Pyrex).
  \item \textbf{Ciments et bétons} : ciment Portland (CaO + SiO$_2$ + Al$_2$O$_3$).
  \item \textbf{Céramiques techniques} (haute performance) : alumine Al$_2$O$_3$, SiC,
        Si$_3$N$_4$, zircone ZrO$_2$\ldots{} plus fibres de renfort, revêtements (diamant),
        abrasifs.
\end{enumerate}
\end{defnbox}

Elles sont faites des éléments les plus \textbf{abondants} de l'écorce terrestre (O, Si, Al,
Ca, Na, Mg, K\ldots) : beaucoup sont bon marché.

\subsection*{Choisir un matériau léger et rigide : les indices $E/\rho$ et $E/\rho^2$}

Les cartes d'Ashby placent chaque matériau dans un plan (densité, module\ldots) en échelle
logarithmique. Pour s'en servir, il faut savoir \textbf{quelle combinaison de propriétés} on
veut maximiser, et cela dépend de la forme de la pièce.

\begin{itemize}
  \item \textbf{Barreau en traction} (rayon $r$, longueur $L$, flèche $\delta$ sous la force $F$) :
        $\dfrac{F}{\pi r^2} = E\dfrac{\delta}{L}$ et $M = \pi r^2 L\rho$, d'où
        $M = \dfrac{L^2 F}{\delta}\,\dfrac{\rho}{E}$. Pour la masse la plus faible à rigidité
        donnée : \textbf{maximiser $E/\rho$}.
  \item \textbf{Poutre pleine en flexion} (section carrée $t\times t$, longueur $l$) :
        $\delta = \dfrac{4l^3F}{Et^4}$ et $M = lt^2\rho$, d'où
        $M = \Big(\dfrac{4l^5F}{\delta}\Big)^{1/2}\dfrac{\rho}{E^{1/2}}$ :
        \textbf{maximiser $E^{1/2}/\rho$}, c'est-à-dire $E/\rho^2$.
  \item \textbf{Tube mince en flexion} (rayon $r$ fixé, épaisseur $t \ll r$) : seule $t$ est
        libre, $M \propto \rho/E$ : de nouveau \textbf{$E/\rho$}.
\end{itemize}

Pourquoi la flexion d'une poutre pleine donne-t-elle $\rho^2$ ? Parce qu'en flexion la rigidité
croît comme $t^4$ alors que la masse croît comme $t^2$ : épaissir une poutre légère est très
efficace, donc la densité \og pèse\fg{} deux fois plus.

\slide[0.6]{c02-21}{Carte module--densité : les droites $E/\rho = $ cste et $E/\rho^2 = $ cste. Les céramiques techniques sont en haut : très rigides pour leur poids.}

\begin{propbox}[title=Ce que disent les cartes]
\begin{itemize}
  \item \textbf{Module} : très élevé (alumine 380\,GPa, SiC 410\,GPa, diamant 1050\,GPa).
  \item \textbf{Densité} moyenne (2,2 à 4, sauf zircone et carbures lourds) : excellentes en
        $E/\rho$.
  \item \textbf{Limite d'élasticité} énorme, mais \textbf{sans importance} : elles cassent
        avant de l'atteindre.
  \item \textbf{Réfractaires} (points de fusion élevés) et souvent résistantes à l'attaque
        chimique.
  \item \textbf{Ténacité très faible} ($K_c \approx 0{,}7$ à $\approx 10\,\text{MPa}\sqrt{\text m}$,
        contre 20 à 200 pour les métaux) : \textbf{c'est leur grande faiblesse}.
\end{itemize}
\end{propbox}

\slide[0.6]{c02-26}{Carte ténacité--module : les céramiques sont à droite (rigides) mais en bas (fragiles).}

\begin{goatbox}
Les céramiques, c'est la pierre et le verre : très rigides, très dures, légères pour leur
rigidité, elles supportent la chaleur et la rouille ne les touche pas. Mais elles ont un gros
défaut : la moindre fissure les fait casser net. Pour choisir un matériau, on ne regarde pas
une seule propriété mais un mélange (rigidité divisée par poids, par exemple), et ce mélange
dépend de la forme de la pièce.
\end{goatbox}
